\documentclass[12pt,a4paper]{article}

\usepackage[utf8]{inputenc}
\usepackage[vietnamese]{babel}
\usepackage{amsmath,amssymb,amsthm}
\usepackage{geometry}
\usepackage{enumitem}
\usepackage[most]{tcolorbox}
\usepackage{hyperref}
\usepackage{fancyhdr}
\usepackage{tikz}
\usepackage{eso-pic}
\usepackage{xcolor}

\geometry{a4paper, margin=2cm, bottom=2.5cm}
\hypersetup{colorlinks=true, linkcolor=blue!70!black, urlcolor=blue}
\setlist[itemize]{topsep=4pt, itemsep=2pt}
\setlist[enumerate]{topsep=4pt, itemsep=2pt}

\AddToShipoutPictureFG{%
  \AtPageCenter{%
    \begin{tikzpicture}[remember picture, overlay]
      \node[rotate=45, scale=1, opacity=0.06]{\fontsize{60}{72}\selectfont\bfseries\sffamily Đặng Tấn Quí};
    \end{tikzpicture}%
  }%
}

\pagestyle{fancy}
\fancyhf{}
\fancyhead[L]{\small\textit{Xác suất và Thống kê}}
\fancyhead[R]{\small\textit{Đặng Tấn Quí}}
\fancyfoot[C]{\thepage}
\fancyfoot[L]{\footnotesize\textcopyright\ Đặng Tấn Quí}
\fancyfoot[R]{\footnotesize SĐT: 0377\,122\,210}
\renewcommand{\headrulewidth}{0.4pt}
\renewcommand{\footrulewidth}{0.4pt}
\fancypagestyle{plain}{\fancyhf{}\fancyfoot[C]{\thepage}\fancyfoot[L]{\footnotesize\textcopyright\ Đặng Tấn Quí}\fancyfoot[R]{\footnotesize SĐT: 0377\,122\,210}\renewcommand{\headrulewidth}{0pt}\renewcommand{\footrulewidth}{0.4pt}}

\newtcolorbox{dinhnghia}[1][]{enhanced, breakable, colback=blue!3!white, colframe=blue!60!black, fonttitle=\bfseries, title={Định nghĩa}, #1}
\newtcolorbox{dinhly}[1][]{enhanced, breakable, colback=green!3!white, colframe=green!50!black, fonttitle=\bfseries, title={Định lý}, #1}
\newtcolorbox{tinhchat}[1][]{enhanced, breakable, colback=yellow!5!white, colframe=orange!50!black, fonttitle=\bfseries, title={Tính chất}, #1}
\newtcolorbox{phuongphap}[1][]{enhanced, breakable, colback=gray!5!white, colframe=gray!50!black, fonttitle=\bfseries, title={Phương pháp}, #1}
\newtcolorbox{luuy}[1][]{enhanced, breakable, colback=red!3!white, colframe=red!50!black, fonttitle=\bfseries, title={Lưu ý}, #1}
\newtcolorbox{congthuc}[1][]{enhanced, breakable, colback=cyan!3!white, colframe=cyan!50!black, fonttitle=\bfseries, title={Công thức}, #1}
\newtcolorbox{vidu}[1][]{enhanced, breakable, colback=violet!3!white, colframe=violet!50!black, fonttitle=\bfseries, title={Ví dụ}, #1}

\begin{document}

\thispagestyle{empty}
\begin{tikzpicture}[remember picture, overlay]
  \fill[blue!80!black] (current page.south west) rectangle (current page.north east);
  \fill[blue!60!cyan, opacity=0.8] (current page.south west) -- (current page.north west) -- ([xshift=-3cm]current page.north east) -- (current page.south east) -- cycle;
  \foreach \r/\o in {6/0.05, 4.5/0.08, 3/0.12} {\fill[white, opacity=\o] ([xshift=2cm, yshift=-2cm]current page.north east) circle (\r cm);}
  \draw[white, line width=2pt] ([yshift=-5cm]current page.north west) ++(2cm,0) -- ++(17cm,0);
  \draw[white, line width=2pt] ([yshift=-16cm]current page.north west) ++(2cm,0) -- ++(17cm,0);
  \node[anchor=west, white, font=\fontsize{26}{32}\bfseries\selectfont] at ([xshift=3cm, yshift=-7.5cm]current page.north west) {PHẦN 11: XÁC SUẤT};
  \node[anchor=west, white, font=\fontsize{22}{28}\bfseries\selectfont] at ([xshift=3cm, yshift=-9.2cm]current page.north west) {VÀ THỐNG KÊ};
  \node[anchor=west, white, font=\fontsize{16}{20}\selectfont] at ([xshift=3cm, yshift=-11cm]current page.north west) {Tài liệu luyện thi du học};
  \node[white, opacity=0.10, font=\fontsize{80}{80}\selectfont, rotate=-15] at ([xshift=-3cm, yshift=4cm]current page.center) {$\sigma$};
  \node[anchor=west, white, font=\Large\bfseries] at ([xshift=3cm, yshift=-13cm]current page.north west) {Biên soạn:};
  \node[anchor=west, yellow!80!white, font=\fontsize{28}{34}\bfseries\selectfont] at ([xshift=3cm, yshift=-14.5cm]current page.north west) {ĐẶNG TẤN QUÍ};
  \node[anchor=west, white!90, font=\large] at ([xshift=3cm, yshift=-18cm]current page.north west) {SĐT: 0377\,122\,210};
  \node[anchor=south, white!80, font=\small] at ([yshift=1.5cm]current page.south) {Tài liệu có bản quyền --- Cấm sao chép khi chưa được phép};
  \node[anchor=south east, white, font=\Large\bfseries] at ([xshift=-2cm, yshift=3cm]current page.south east) {\the\year};
\end{tikzpicture}
\newpage

\thispagestyle{empty}
\vspace*{5cm}
\begin{center}
  {\Large\bfseries PHẦN 11: XÁC SUẤT VÀ THỐNG KÊ}\\[8pt]
  {\large Tài liệu luyện thi du học}
\end{center}
\vspace{2cm}
\begin{tcolorbox}[enhanced, colback=red!3!white, colframe=red!60!black, fonttitle=\bfseries, title={Thông tin bản quyền}, boxrule=1.5pt]
  \textbf{Tác giả:} Đặng Tấn Quí \quad \textbf{Liên hệ:} 0377\,122\,210 \quad \textbf{Năm:} \the\year \\[6pt]
  \textbf{Bản quyền \textcopyright\ \the\year\ Đặng Tấn Quí. Bảo lưu mọi quyền.}
\end{tcolorbox}
\vfill\begin{center}\small\textit{Tài liệu lưu hành nội bộ}\end{center}
\newpage
\tableofcontents
\newpage

%% ================================================================
%%  PHẦN 11: XÁC SUẤT VÀ THỐNG KÊ
%% ================================================================
\section{Xác suất và Thống kê}

%% ----------------------------------------------------------------
%%  11.1 Các loại phân phối
%% ----------------------------------------------------------------
\subsection{Các loại phân phối: Phân phối Nhị thức, Poisson và Chuẩn}

\subsubsection{Phân phối Nhị thức (Binomial Distribution)}

\begin{dinhnghia}[title={Thí nghiệm Bernoulli và phân phối nhị thức}]
  \textbf{Thí nghiệm Bernoulli} là thí nghiệm chỉ có hai kết quả: ``thành công'' (xác suất $p$) và ``thất bại'' (xác suất $1-p = q$).

  Nếu thực hiện $n$ thí nghiệm Bernoulli độc lập với cùng xác suất thành công $p$, biến ngẫu nhiên $X$ = ``số lần thành công'' có \textbf{phân phối nhị thức}: $X \sim B(n, p)$.
\end{dinhnghia}

\begin{congthuc}[title={Phân phối nhị thức}]
  \[
    P(X = k) = \binom{n}{k} p^k (1-p)^{n-k}, \quad k = 0, 1, 2, \ldots, n.
  \]
  Với $\binom{n}{k} = \dfrac{n!}{k!(n-k)!}$ là hệ số tổ hợp.

  \textbf{Kỳ vọng và phương sai:}
  \[
    \mu = E[X] = np, \qquad \sigma^2 = \text{Var}(X) = np(1-p) = npq.
  \]
\end{congthuc}

\begin{vidu}
  Gieo đồng xu (xác suất sấp $p=0.5$) $n=10$ lần. Tính xác suất được đúng $4$ lần sấp.

  \medskip\textbf{Giải:}
  $X \sim B(10, 0.5)$.
  $P(X=4) = \binom{10}{4}(0.5)^4(0.5)^6 = 210 \times (0.5)^{10} = \dfrac{210}{1024} \approx 0.205$.
\end{vidu}

\subsubsection{Phân phối Poisson}

\begin{dinhnghia}[title={Phân phối Poisson}]
  Biến ngẫu nhiên $X$ có \textbf{phân phối Poisson} với tham số $\lambda > 0$ ($\lambda$ là tốc độ trung bình xảy ra sự kiện) nếu:
  \[
    P(X = k) = \frac{e^{-\lambda} \lambda^k}{k!}, \quad k = 0, 1, 2, \ldots
  \]
  \textbf{Kỳ vọng và phương sai:}
  \[
    \mu = E[X] = \lambda, \qquad \sigma^2 = \text{Var}(X) = \lambda.
  \]
\end{dinhnghia}

\begin{luuy}
  Phân phối Poisson thường dùng để mô hình hóa:
  \begin{itemize}
    \item Số cuộc gọi đến tổng đài trong một giờ.
    \item Số lỗi trên $1\,\text{m}$ sợi dây.
    \item Số khách hàng đến trong một khoảng thời gian.
  \end{itemize}
  \textbf{Xấp xỉ nhị thức:} Khi $n$ lớn, $p$ nhỏ và $\lambda = np$ (trung bình): $B(n,p) \approx \text{Poisson}(\lambda)$.
\end{luuy}

\begin{vidu}
  Trung bình có $\lambda = 3$ xe đến trạm xăng mỗi phút. Tính xác suất đúng $5$ xe đến trong một phút.

  \medskip\textbf{Giải:}
  $P(X=5) = \dfrac{e^{-3} \cdot 3^5}{5!} = \dfrac{e^{-3} \cdot 243}{120} \approx \dfrac{0.04979 \times 243}{120} \approx 0.1008$.
\end{vidu}

\subsubsection{Phân phối Chuẩn (Normal Distribution)}

\begin{dinhnghia}[title={Phân phối chuẩn}]
  Biến ngẫu nhiên liên tục $X$ có \textbf{phân phối chuẩn} với tham số $\mu$ (kỳ vọng) và $\sigma^2$ (phương sai), ký hiệu $X \sim N(\mu, \sigma^2)$, nếu hàm mật độ xác suất là:
  \[
    f(x) = \frac{1}{\sigma\sqrt{2\pi}} e^{-\frac{(x-\mu)^2}{2\sigma^2}}, \quad x \in \mathbb{R}.
  \]
\end{dinhnghia}

\begin{dinhnghia}[title={Phân phối chuẩn tắc (Standard Normal)}]
  \textbf{Phân phối chuẩn tắc} $Z \sim N(0,1)$ có hàm mật độ:
  \[
    \phi(z) = \frac{1}{\sqrt{2\pi}} e^{-z^2/2}.
  \]
  và hàm phân phối tích lũy: $\Phi(z) = P(Z \le z) = \displaystyle\int_{-\infty}^z \phi(t)\,dt$.
\end{dinhnghia}

\begin{phuongphap}[title={Chuẩn hóa (Standardization)}]
  Nếu $X \sim N(\mu, \sigma^2)$, đặt $Z = \dfrac{X - \mu}{\sigma}$, thì $Z \sim N(0,1)$.

  \[
    P(a \le X \le b) = P\!\left(\frac{a-\mu}{\sigma} \le Z \le \frac{b-\mu}{\sigma}\right) = \Phi\!\left(\frac{b-\mu}{\sigma}\right) - \Phi\!\left(\frac{a-\mu}{\sigma}\right).
  \]
\end{phuongphap}

\begin{tinhchat}[title={Quy tắc Empirical (68--95--99.7)}]
  Với $X \sim N(\mu, \sigma^2)$:
  \begin{itemize}
    \item $P(\mu - \sigma < X < \mu + \sigma) \approx 68\%$.
    \item $P(\mu - 2\sigma < X < \mu + 2\sigma) \approx 95\%$.
    \item $P(\mu - 3\sigma < X < \mu + 3\sigma) \approx 99.7\%$.
  \end{itemize}
\end{tinhchat}

\begin{vidu}
  Chiều cao học sinh $X \sim N(165, 36)$ cm (tức $\mu = 165$, $\sigma = 6$). Tìm $P(159 < X < 177)$.

  \medskip\textbf{Giải:}
  $P(159 < X < 177) = P\!\left(\dfrac{159-165}{6} < Z < \dfrac{177-165}{6}\right) = P(-1 < Z < 2) = \Phi(2) - \Phi(-1)$.

  $= \Phi(2) - (1 - \Phi(1)) \approx 0.9772 - (1 - 0.8413) = 0.9772 - 0.1587 = 0.8185$.
\end{vidu}

%% ----------------------------------------------------------------
%%  11.2 Kiểm định giả thuyết
%% ----------------------------------------------------------------
\subsection{Kiểm định giả thuyết (Hypothesis Testing)}

\begin{dinhnghia}[title={Khái niệm cơ bản}]
  \begin{itemize}
    \item \textbf{Giả thuyết không ($H_0$):} Phát biểu mặc định, ``không có hiệu ứng'', cần kiểm tra.
    \item \textbf{Giả thuyết thay thế ($H_1$ hay $H_a$):} Phát biểu ta muốn ủng hộ.
    \item \textbf{Mức ý nghĩa ($\alpha$):} Xác suất mắc \textbf{lỗi loại I} (bác bỏ $H_0$ đúng). Thường $\alpha = 0.05$ hoặc $0.01$.
    \item \textbf{Giá trị p (p-value):} Xác suất thu được kết quả mẫu cực đoan như quan sát (hoặc cực đoan hơn) nếu $H_0$ đúng.
    \item \textbf{Quy tắc quyết định:} Bác bỏ $H_0$ nếu p-value $\le \alpha$.
  \end{itemize}
\end{dinhnghia}

\begin{congthuc}[title={Kiểm định $z$ (Z-test) -- trung bình tổng thể}]
  Khi mẫu lớn ($n \ge 30$) hoặc biết $\sigma$: thống kê kiểm định:
  \[
    Z = \frac{\bar{X} - \mu_0}{\sigma/\sqrt{n}} \sim N(0,1) \quad \text{(dưới } H_0).
  \]
  \begin{itemize}
    \item $H_0: \mu = \mu_0$ so với $H_1: \mu \ne \mu_0$ (hai phía): bác bỏ nếu $|Z| > z_{\alpha/2}$.
    \item $H_1: \mu > \mu_0$ (một phía phải): bác bỏ nếu $Z > z_\alpha$.
    \item $H_1: \mu < \mu_0$ (một phía trái): bác bỏ nếu $Z < -z_\alpha$.
  \end{itemize}
\end{congthuc}

\begin{congthuc}[title={Kiểm định $t$ (T-test) -- mẫu nhỏ}]
  Khi mẫu nhỏ ($n < 30$) và không biết $\sigma$: thống kê kiểm định:
  \[
    T = \frac{\bar{X} - \mu_0}{s/\sqrt{n}} \sim t_{n-1} \quad \text{(phân phối t với } n-1 \text{ bậc tự do)}.
  \]
  trong đó $s$ là độ lệch chuẩn mẫu: $s = \sqrt{\dfrac{\sum(x_i - \bar{x})^2}{n-1}}$.
\end{congthuc}

\begin{congthuc}[title={Kiểm định Chi-squared ($\chi^2$)}]
  Kiểm định sự phù hợp (goodness-of-fit) hoặc độc lập trong bảng phân phối tần số:
  \[
    \chi^2 = \sum \frac{(O_i - E_i)^2}{E_i},
  \]
  trong đó $O_i$ là tần số quan sát và $E_i$ là tần số kỳ vọng. So sánh với $\chi^2_{k-1}$ ($k$ là số hạng).
\end{congthuc}

\begin{phuongphap}[title={Quy trình kiểm định giả thuyết}]
  \begin{enumerate}
    \item Phát biểu $H_0$ và $H_1$; chọn mức ý nghĩa $\alpha$.
    \item Chọn phép kiểm định phù hợp ($z$-test, $t$-test, $\chi^2$-test, \ldots).
    \item Tính thống kê kiểm định từ dữ liệu mẫu.
    \item Tính p-value (hoặc so sánh với giá trị tới hạn).
    \item Ra quyết định: bác bỏ hay không bác bỏ $H_0$.
    \item Kết luận trong bối cảnh bài toán.
  \end{enumerate}
\end{phuongphap}

\begin{vidu}
  Một nhà sản xuất tuyên bố trọng lượng trung bình bánh $\mu_0 = 200\,\text{g}$. Mẫu $n=36$ bánh cho $\bar{x} = 197.5\,\text{g}$, $\sigma = 6\,\text{g}$. Kiểm định với $\alpha = 0.05$ có đủ bằng chứng $\mu < 200$ không?

  \medskip\textbf{Giải:}
  $H_0: \mu = 200$, $H_1: \mu < 200$ (một phía trái).

  $Z = \dfrac{197.5 - 200}{6/\sqrt{36}} = \dfrac{-2.5}{1} = -2.5$.

  $z_\alpha = z_{0.05} = 1.645$ (một phía). Bác bỏ $H_0$ nếu $Z < -1.645$.

  $Z = -2.5 < -1.645$: \textbf{bác bỏ $H_0$}. Có đủ bằng chứng trọng lượng thực nhỏ hơn $200\,\text{g}$.
\end{vidu}

%% ----------------------------------------------------------------
%%  11.3 Hàm mật độ xác suất (PDF)
%% ----------------------------------------------------------------
\subsection{Hàm mật độ xác suất (PDF): Kết hợp giải tích để tính xác suất}

\begin{dinhnghia}[title={Biến ngẫu nhiên liên tục và PDF}]
  Biến ngẫu nhiên liên tục $X$ có \textbf{hàm mật độ xác suất (PDF)} $f(x)$ nếu:
  \begin{enumerate}
    \item $f(x) \ge 0$ với mọi $x$.
    \item $\displaystyle\int_{-\infty}^{+\infty} f(x)\,dx = 1$.
    \item $P(a \le X \le b) = \displaystyle\int_a^b f(x)\,dx$.
  \end{enumerate}
  \textbf{Hàm phân phối tích lũy (CDF):} $F(x) = P(X \le x) = \displaystyle\int_{-\infty}^x f(t)\,dt$, và $F'(x) = f(x)$.
\end{dinhnghia}

\begin{congthuc}[title={Kỳ vọng và phương sai của biến liên tục}]
  \begin{itemize}
    \item \textbf{Kỳ vọng:} $E[X] = \displaystyle\int_{-\infty}^{+\infty} x\,f(x)\,dx$.
    \item \textbf{Phương sai:} $\text{Var}(X) = E[(X-\mu)^2] = \displaystyle\int_{-\infty}^{+\infty} (x-\mu)^2 f(x)\,dx = E[X^2] - (E[X])^2$.
    \item \textbf{Kỳ vọng hàm $g(X)$:} $E[g(X)] = \displaystyle\int_{-\infty}^{+\infty} g(x)\,f(x)\,dx$.
  \end{itemize}
\end{congthuc}

\begin{vidu}
  Cho PDF $f(x) = \begin{cases} cx^2 & 0 \le x \le 3 \\ 0 & \text{khác}\end{cases}$.
  \begin{enumerate}[label=(\alph*)]
    \item Tìm $c$.
    \item Tính $P(1 \le X \le 2)$ và $E[X]$.
  \end{enumerate}

  \medskip\textbf{Giải:}
  \begin{enumerate}[label=(\alph*)]
    \item $\displaystyle\int_0^3 cx^2\,dx = c\cdot\dfrac{27}{3} = 9c = 1 \Rightarrow c = \dfrac{1}{9}$.

    \item $P(1 \le X \le 2) = \displaystyle\int_1^2 \dfrac{x^2}{9}\,dx = \dfrac{1}{9}\cdot\dfrac{x^3}{3}\Big|_1^2 = \dfrac{1}{27}(8-1) = \dfrac{7}{27}$.

      $E[X] = \displaystyle\int_0^3 x \cdot \dfrac{x^2}{9}\,dx = \dfrac{1}{9}\int_0^3 x^3\,dx = \dfrac{1}{9}\cdot\dfrac{81}{4} = \dfrac{9}{4}$.
  \end{enumerate}
\end{vidu}

\begin{congthuc}[title={Phân phối đều liên tục}]
  $X \sim U(a,b)$: PDF đều trên $[a,b]$:
  \[
    f(x) = \frac{1}{b-a}, \quad a \le x \le b.
  \]
  $E[X] = \dfrac{a+b}{2}$, $\quad \text{Var}(X) = \dfrac{(b-a)^2}{12}$.
\end{congthuc}

\begin{congthuc}[title={Phân phối mũ (Exponential)}]
  $X \sim \text{Exp}(\lambda)$: PDF $f(x) = \lambda e^{-\lambda x}$ với $x \ge 0$.
  \[
    E[X] = \frac{1}{\lambda}, \qquad \text{Var}(X) = \frac{1}{\lambda^2}.
  \]
  \textbf{Tính không có ký ức:} $P(X > s+t \mid X > s) = P(X > t)$.
\end{congthuc}

\medskip\noindent\textbf{Các dạng bài tập tổng hợp:}
\begin{enumerate}[label=\textbf{Dạng \arabic*:}, leftmargin=*, align=left]
  \item Tính xác suất dùng phân phối nhị thức.
  \item Tính xác suất dùng phân phối Poisson; xấp xỉ nhị thức bằng Poisson.
  \item Chuẩn hóa và tính xác suất dùng phân phối chuẩn (tra bảng hoặc tính).
  \item Kiểm định $z$-test (trung bình với $\sigma$ biết).
  \item Kiểm định $t$-test (trung bình với $\sigma$ không biết, mẫu nhỏ).
  \item Kiểm định $\chi^2$ (bảng tần số, độc lập hay phù hợp).
  \item Tìm hằng số $c$ để hàm cho trước là PDF; tính xác suất, kỳ vọng, phương sai.
  \item Tìm CDF từ PDF; tìm trung vị, phân vị.
\end{enumerate}

\end{document}
